\begin{ouverture}[college3]{14}{Mesurer l’intensité du courant}{Partie 3 · L’énergie, ses transferts et ses conversions}{Le courant est-il consommé par une lampe ?}
\objectif{Brancher un ampèremètre en série et lire son calibre.}
\objectif{Appliquer l’unicité en série et l’additivité aux nœuds.}
\objectif{Argumenter à partir de mesures et de leur précision.}
\prerequis{Circuits en série et en dérivation.}
\end{ouverture}
\pcPage{Le cours}{Mesurer l’intensité du courant}
\phantomsection\label{ess-14}
\essSection{1}{Une grandeur, une unité}
L’intensité $I$ caractérise le courant électrique ; son unité est l’ampère (A). On utilise souvent le milliampère : $1\unite A=1\,000\unite{mA}$. Une lampe ne « consomme » pas le courant : elle reçoit de l’énergie transférée électriquement et la convertit. L’intensité et l’énergie ne sont pas la même grandeur.\par
\pcFigure{figures/c14-amperemetre.pdf}{L’ampèremètre est inséré dans la branche ; les charges doivent traverser l’appareil, qui ne se branche jamais directement aux bornes du générateur.}
\essSection{2}{Insérer l’ampèremètre}
Ouvrir le circuit hors alimentation, puis insérer l’appareil en série. Choisir la fonction courant continu et les bornes prescrites ; commencer avec un calibre adapté élevé puis affiner si possible. Le courant conventionnel entre par la borne A ou mA et sort par COM pour une lecture positive. Un signe négatif indique le sens de branchement dans le modèle, pas une intensité physiquement « impossible ».\par
\clearpage
\essSection{3}{Une même intensité en série}
Dans un circuit à une seule branche de dipôles en série, l’intensité est la même en tout point en régime établi. Pour deux lampes : $I_1=I_2=I_G$. Déplacer correctement l’ampèremètre permet de tester cette loi. Ajouter un dipôle peut modifier la valeur du courant dans toute la boucle ; cela ne contredit pas l’unicité à un instant donné.\par
\pcFigure{figures/c14-noeud.pdf}{Au nœud, la somme des intensités entrantes égale celle des intensités sortantes. Les flèches fixent les sens choisis.}
\essSection{4}{L’additivité en dérivation}
Au nœud de séparation, dans les sens du schéma, $I_G=I_1+I_2$. Les intensités des branches ne sont pas forcément égales : elles dépendent des dipôles. Pour vérifier, mesurer les trois valeurs sans modifier le fonctionnement du circuit. Comparer les sommes en tenant compte de la résolution de l’appareil et des variations de la source.\par
\begin{pcEncadre}{method}{Un exemple expliqué}Un générateur fournit $I_G=0{,}45\unite A$ à deux branches. On mesure $I_1=180\unite{mA}=0{,}180\unite A$. Alors $I_2=I_G-I_1=0{,}45-0{,}180=0{,}270\unite A=270\unite{mA}$. Les unités sont harmonisées avant la soustraction.\end{pcEncadre}
\pcPage{Expérimenter et retenir}{Mesurer l’intensité du courant}
\phantomsection\label{activite-14}
\begin{pcEncadre}{method}{Investigation en classe}\textbf{Question.} Le courant diminue-t-il après une lampe ? Dans un circuit scolaire série, mesurer successivement l’intensité avant et après une lampe, en interrompant l’alimentation pour modifier le montage. Recommencer avec deux lampes et comparer. Ensuite, sous conduite du professeur, mesurer une branche principale et deux branches dérivées. \textbf{À rendre :} schémas, calibres, tableau des mesures, somme des branches et conclusion. L’ampèremètre n’est jamais branché directement en dérivation sur une source.\end{pcEncadre}
\pcSousTitre{Ce que je dois savoir expliquer}\begin{questions}
\item L’intensité se mesure en A, avec un ampèremètre en série.
\item 1 A correspond à 1 000 mA.
\item L’intensité est identique en tout point d’une branche série en régime établi.
\item Au nœud, les intensités s’additionnent avec les sens fixés.
\item La lampe convertit de l’énergie ; elle ne consomme pas le courant.
\end{questions}
\begin{pcEncadre}{attention}{Une erreur à éviter}Brancher un ampèremètre comme un voltmètre peut court-circuiter la source et détruire le fusible de l’appareil ; vérifier le schéma hors alimentation.\end{pcEncadre}
\pcRepere{cor-activite-14}{Pistes pour l’investigation}\pcGrille
\clearpage\markright{Je démarre}
\rubrique{Je démarre}
\par\noindent\begin{minipage}{\linewidth}
\begin{exercice}[1]{L’unité}
\begin{questions}
\item Quelle unité et quel symbole pour l’intensité ?
\item Quel instrument la mesure ?
\end{questions}
\end{exercice}
\end{minipage}
\par\noindent\begin{minipage}{\linewidth}
\begin{exercice}[1]{Le branchement}
\begin{questions}
\item Un ampèremètre se place-t-il en série ou en dérivation ?
\item Faut-il couper l’alimentation pour modifier le circuit ?
\end{questions}
\end{exercice}
\end{minipage}
\par\noindent\begin{minipage}{\linewidth}
\begin{exercice}[1]{Convertir}
\begin{questions}
\item Convertir $250\unite{mA}$ en A.
\item Convertir $0{,}08\unite A$ en mA.
\end{questions}
\end{exercice}
\end{minipage}
\par\noindent\begin{minipage}{\linewidth}
\begin{exercice}[1]{Dans une série}
\begin{questions}
\item Avant une lampe, $I=0{,}12\unite A$. Quelle intensité après, dans la même branche ?
\item La lampe a-t-elle consommé une partie du courant ?
\end{questions}
\end{exercice}
\end{minipage}
\par\noindent\begin{minipage}{\linewidth}
\begin{exercice}[1]{Un nœud simple}
\pcFigureEx{figures/ex-noeud.pdf}{}
\begin{questions}
\item Deux courants de $0{,}10$ et $0{,}20\unite A$ se réunissent. Quelle intensité sortante ?
\item Quelle loi utiliser ?
\end{questions}
\end{exercice}
\end{minipage}
\clearpage\markright{Je m’entraîne}
\rubrique{Je m’entraîne}
\par\noindent\begin{minipage}{\linewidth}
\begin{exercice}[2]{Une branche inconnue}
\pcFigureEx{figures/ex-courants-050.pdf}{}
\begin{questions}
\item $I_G=0{,}50\unite A$ et $I_1=0{,}18\unite A$. Calculer $I_2$.
\item Vérifier la somme.
\end{questions}
\end{exercice}
\end{minipage}
\par\noindent\begin{minipage}{\linewidth}
\begin{exercice}[2]{Attention aux unités}
\begin{questions}
\item $I_G=600\unite{mA}$ et $I_1=0{,}25\unite A$. Trouver $I_2$ en mA.
\item Puis en A.
\end{questions}
\end{exercice}
\end{minipage}
\par\noindent\begin{minipage}{\linewidth}
\begin{exercice}[2]{Choisir un calibre}
\begin{questions}
\item On attend environ $150\unite{mA}$. Entre les calibres 20 mA, 200 mA et 10 A, lequel permet une mesure précise sans dépassement attendu, une fois le courant vérifié ?
\item Pourquoi commencer prudemment ?
\end{questions}
\end{exercice}
\end{minipage}
\par\noindent\begin{minipage}{\linewidth}
\begin{exercice}[2]{Le signe négatif}
\begin{questions}
\item Un multimètre affiche $-0{,}15\unite A$. Que peut-on dire du sens du courant par rapport au branchement choisi ?
\item Quelle valeur de l’intensité dans ce modèle ?
\end{questions}
\end{exercice}
\end{minipage}
\par\noindent\begin{minipage}{\linewidth}
\begin{exercice}[2]{Deux branches identiques}
\begin{questions}
\item Deux branches strictement identiques sont alimentées de la même façon et le courant principal vaut $0{,}40\unite A$. Quel courant par branche dans le modèle symétrique ?
\item La division par deux vaut-elle pour toutes branches ?
\end{questions}
\end{exercice}
\end{minipage}
\clearpage\markright{Je raisonne}
\rubrique{Je raisonne}
\par\noindent\begin{minipage}{\linewidth}
\begin{exercice}[3]{Une expérience cohérente ?}
\begin{questions}
\item On mesure $I_G=0{,}300\unite A$, $I_1=0{,}122\unite A$ et $I_2=0{,}176\unite A$. Calculer l’écart de bilan.
\item Faut-il rejeter la loi immédiatement ?
\end{questions}
\end{exercice}
\end{minipage}
\par\noindent\begin{minipage}{\linewidth}
\begin{exercice}[3]{Un emplacement dangereux}
\pcFigureEx{figures/ex-amperemetre-danger.pdf}{Montage incorrect à analyser ; ne pas le réaliser.}
\begin{questions}
\item Un élève branche un ampèremètre directement entre les deux bornes d’une pile. Pourquoi le montage est-il incorrect ?
\item Quelle correction conceptuelle proposer ?
\end{questions}
\end{exercice}
\end{minipage}
\par\noindent\begin{minipage}{\linewidth}
\begin{exercice}[3]{Une lampe supplémentaire}
\begin{questions}
\item On ajoute une lampe en série et le courant passe de 0,20 à 0,12 A dans toute la boucle. Est-ce contraire à l’unicité ?
\item Que compare la loi ?
\end{questions}
\end{exercice}
\end{minipage}
\par\noindent\begin{minipage}{\linewidth}
\begin{exercice}[3]{Une branche ouverte}
\begin{questions}
\item Dans un circuit, $I_1=0{,}15\unite A$ et $I_2=0{,}10\unite A$. On ouvre la branche 2. Quelle relation subsiste ?
\item Peut-on affirmer que $I_1$ reste exactement identique avec toute source réelle ?
\end{questions}
\end{exercice}
\end{minipage}
\par\noindent\begin{minipage}{\linewidth}
\begin{exercice}[3]{Trois branches}
\pcFigureEx{figures/ex-trois-courants.pdf}{}
\begin{questions}
\item Le courant principal vaut $0{,}80\unite A$ ; deux branches valent $220\unite{mA}$ et $0{,}35\unite A$. Calculer la troisième.
\item Donner un contrôle rapide.
\end{questions}
\end{exercice}
\end{minipage}
\pcPage{Défi du chapitre}{Mesurer l’intensité du courant}
\phantomsection\label{defi-ex-14}
\pcDefiTitre{14}{Trouver la mesure incohérente}
\pcFigureEx{figures/ex-bilan-intensites.pdf}{}
\begin{questions}
\item Un relevé indique $I_G=0{,}40\unite A$, $I_1=260\unite{mA}$ et $I_2=0{,}24\unite A$. Vérifier le bilan.
\item Si les deux branches sont fiables, quel courant principal attendre ?
\item Proposer trois vérifications avant de conclure à une loi fausse.
\end{questions}
\pcGrille
